% \iffalse
%<*driver>
\ProvidesFile{seqenseigne-carte-automatisme.dtx}%
[2026/05/12 v0.13.6.2 seqenseigne — Module cartes d'automatisme (Leitner)]
%</driver>
% \fi
%
% =============================================================================
% seqenseigne-carte-automatisme.dtx
% Module : CARTE-AUTOMATISME
% Rôle   :
%   - Définir le format 48×69mm pour les cartes d'automatisme
%   - Fournir les macros
%		- seqInitCarteParam pour inialiser la production 
%			de planches de cartes paramétrées
%		- \seqCarteRecto pour produire 1 recto
%		- \seqCarteVerso pour produire 1 verso
%		- \seqCarteAuto pour produire 1 carte complète (recto/verso)
%		- \seqCartePageRecto pour produire 1 page de rectos identiques
%		- \seqCartePageVerso pour produire 1 page de versos identiques
%     
%
% Histoire :
%   Ce module REMPLACE l'ancien seqenseigne-flashcard.dtx (v1.0-dev 2025/03,
%   resté à l'état d'expérimentation et jamais réellement utilisé). Le
%   modèle pédagogique a évolué : on ne produit plus des cartes par
%   objectif avec critères TB/S/F, mais des cartes d'automatisme
%   recto = question / verso = réponse, à utiliser en rituel Leitner
%   (5 niveaux d'enveloppes, évaluation orale mensuelle).
%
% =============================================================================
% IMPORTANT — Architecture
% =============================================================================
% Ce module est volontairement AGNOSTIQUE de la base de données seqenseigne.
% Il ne lit AUCUN fichier .dbtex (param_niveaux, themes, etc.). Tout ce qui
% est métier (lookup, libellés, couleurs) doit être résolu côté APPLI et
% passé en argument aux macros publiques :
%   - L'identifiant de carte (niveau, sequence, num) : passé tel quel
%   - Le libellé du type pédagogique (`Définition`, `Propriété`, etc.) :
%     passé en clair, pas un mapping local
%   - Le code couleur du thème de la séquence (`nombres`, `geometrie`,
%     etc.) : passé tel quel, le module appelle directement \seq@setColors
%     qui est un dispatcher pur LaTeX (sans dépendance .dbtex)
%
% Conséquence : ce .dtx peut être chargé seul (avec seqenseigne-theme pour
% les couleurs) sans avoir besoin de \seqLoadData ni de fichiers .dbtex.
%
% =============================================================================
% Dépendances :
%   - seqenseigne-theme.dtx (couleurs \seq@couleuri@vif etc. + dispatcher
%     \seq@setColors qui est pur LaTeX)
%   - tcolorbox avec les bibliothèques skins et breakable (chargées ci-dessous)
%
% Macros publiques exposées :
%   \seqCarteAuto[options]{recto}{verso}%
%     options clé=valeur :
%       niveau              : code niveau   (ex. N10) — REQUIS
%       sequence            : code séquence (ex. S01) — REQUIS
%       num                 : numéro carte  (entier 1..99) — REQUIS
%       typepedagolibelle   : libellé en clair affiché dans le bandeau
%                             (ex. 'Définition', 'Propriété', ...) — REQUIS
%                             NB : la clé est sans underscore (xkeyval
%                             ne supporte pas '_' dans les noms de clé).
%       codecouleur         : code symbolique du thème — REQUIS
%                             (l'un de : nombres, donnees, grandeurs,
%                                        geometrie, algorithmique)
%       nom                 : nom interne (optionnel, non affiché)
%
% Usage typique :
%   \documentclass{article}%
%   \usepackage{seqenseigne-core}%
%   \usepackage{seqenseigne-theme}%
%   \usepackage{seqenseigne-carte-automatisme}%
%   \begin{document}%
%     \seqCarteAuto[
%       niveau=N10, sequence=S01, num=3,
%       typepedagolibelle=Propriété,
%       codecouleur=nombres
%     ]
%     {Combien font $20\,\%$ de $100$?}%
%     {$20$}%
%     % Plus de cartes...
%   \end{document}%
%
% L'appli seqenseigne génère ce .tex à partir de la BDD cartes_automatisme.
% =============================================================================
%
%<*package>
\NeedsTeXFormat{LaTeX2e}%
\ProvidesPackage{seqenseigne-carte-automatisme}%
[2026/05/12 v0.13.6.2 seqenseigne carte-automatisme module]

% =============================================================================
% Dépendances
% =============================================================================
\RequirePackage{xkeyval}%
\RequirePackage{xifthen}%
\RequirePackage{etoolbox}%
\RequirePackage{geometry}%
% NB v0.13.6.2.2 : on NE charge PAS tcolorbox ni \tcbuselibrary ici.
% seqenseigne-theme.dtx les charge déjà avec les options
% [theorems,breakable,skins]. Si on les rechargeait ici (même sans option),
% LaTeX produirait "Option clash for package tcolorbox" au moment où
% seqenseigne-theme essaie de re-poser ses options. Le contrat est :
% ce module dépend de seqenseigne-theme, qui apporte tcolorbox avec
% les libraries skins et breakable nécessaires aux clés `enhanced` et
% `attach boxed title to top left` utilisées plus bas.

% =============================================================================
% Géométrie de page
% =============================================================================
%
% \begin{macrocode}%
\geometry{a4paper,landscape,%
	margin=1mm,%
	nohead,nofoot}%
% Pas de numéro de page (les cartes sont des unités autonomes).
\pagestyle{empty}%
% \end{macrocode}%

% =============================================================================
% Couleurs internes du module
% =============================================================================
%
% Couleurs locales du module, initialisées avec des valeurs neutres puis
% remplacées par \seqca@applytheme à chaque appel de \seqCarteAuto.
%
% Ces couleurs sont des MACROS (pas des noms xcolor), donc on les utilise
% comme un argument de colback/colframe (qui acceptent l'expression).
%
\providecommand{\seqca@couleur@titre}{blue!70!black}%
\providecommand{\seqca@couleur@fond}{blue!10!white}%
\providecommand{\seqca@couleur@cadre}{blue!70!black}%
\providecommand{\seqca@couleur@versofond}{white}%

% \begin{macro}{\seqca@applytheme}%
% Synchronise les couleurs internes de la carte avec les couleurs courantes
% du thème seqenseigne (variables \seq@couleuri@vif etc., positionnées par
% \seq@setColors).
% \begin{macrocode}%
\newcommand{\seqca@applytheme}{%
    \renewcommand{\seqca@couleur@titre}{\seq@couleuri@vif}%
    \renewcommand{\seqca@couleur@fond}{\seq@couleurii@pale}%
    \renewcommand{\seqca@couleur@cadre}{\seq@couleuri@vif}%
}%
% \end{macrocode}%
% \end{macro}%

% =============================================================================
% Compteur de position dans la grille 4x4 des planches
% =============================================================================
\newcounter{seqca@cellpos}

% =============================================================================
% MACROS ET ENVIRONNEMENTS INTERNES
% =============================================================================
%
% -----------------------------------------------------------------------------
% Formatage de l'identifiant N10/S01/CA03
% -----------------------------------------------------------------------------
% \begin{macro}{\seqca@identifiant}%
% #1 : code niveau (ex. N10)
% #2 : code séquence (ex. S01)
% #3 : numéro carte (entier, ex. 3)
% Renvoie : N10/S01/CA03 (num padding 2 chiffres)
% \begin{macrocode}%
\newcommand{\seqca@identifiant}[3]{%
    #1/#2/CA\ifnum#3<10 0\fi\number#3%
}%
% \end{macrocode}%
% \end{macro}%
% -----------------------------------------------------------------------------
% Composition du recto
% -----------------------------------------------------------------------------
% \begin{macro}{\seqca@recto}%
% #1 : code niveau
% #2 : code séquence
% #3 : numéro carte
% #4 : libellé en clair du type pédagogique
% #5 : contenu LaTeX du recto (question)
% \begin{macrocode}%
\newcommand{\seqca@recto}[5]{%
    \noindent
    \begin{tcolorbox}[
        width           = 69mm,
        height          = 48mm,
        colback         = \seqca@couleur@fond,
        colframe        = \seqca@couleur@cadre,
        colbacktitle    = \seqca@couleur@titre,
        coltitle        = white,
        fonttitle       = \bfseries\scriptsize,
        title           = {%
            \seqca@identifiant{#1}{#2}{#3} \hfill #4%
        },
        enhanced,
        nobeforeafter,
        boxsep          = 1pt,
        top             = 3pt,
        bottom          = 3pt,
        left            = 3pt,
        right           = 3pt,
        halign			= center,
        valign			= center,
        space to upper
    ]
        % Zone contenu : centrée verticalement
        #5
        % Zone identification élève en bas (Nom / Classe, pas de date)
        \tcblower
        \scriptsize
        Nom~: \underline{\hspace{2.2cm}}\hfill
        Classe~: \underline{\hspace{1.2cm}}%
    \end{tcolorbox}%
}%
% \end{macrocode}%
% \end{macro}%
% -----------------------------------------------------------------------------
% Composition du verso
% -----------------------------------------------------------------------------
% \begin{macro}{\seqca@verso}%
% #1 : code niveau
% #2 : code séquence
% #3 : numéro carte
% #4 : contenu LaTeX du verso (réponse)
% \begin{macrocode}%
\newcommand{\seqca@verso}[4]{%
    \noindent
    \begin{tcolorbox}[
        width           = 69mm,
        height          = 48mm,
        colback         = \seqca@couleur@versofond,
        colframe        = \seqca@couleur@cadre,
        colbacktitle    = \seqca@couleur@titre,
        coltitle        = white,
        fonttitle       = \bfseries\scriptsize,
        title           = {%
            \seqca@identifiant{#1}{#2}{#3} \hfill Réponse%
        },
        enhanced,
        nobeforeafter,
        boxsep          = 1pt,
        top             = 3pt,
        bottom          = 3pt,
        left            = 3pt,
        right           = 3pt,
        halign			= center,
        valign			= center,
        space to upper
    ]
        #4
    \end{tcolorbox}%
}%
% \end{macrocode}%
% \end{macro}%
% -----------------------------------------------------------------------------
% Ecrit le recto dans le fichier annexe
% -----------------------------------------------------------------------------
% \begin{macro}{\seqca@emettreRecto}%
% #1 : contenu LaTeX du recto
% \begin{macrocode}%
\newcommand{\seqca@emettreRecto}[1]{%
	\stepcounter{seqca@cellpos}%
	\immediate\write\plancherectostream{%
		\unexpanded{\seqCarteRecto[niveau=}\cmdseqca@planche@niveau%
		\unexpanded{,sequence=}\cmdseqca@planche@sequence%
		\unexpanded{,num=}\cmdseqca@planche@num%
		\unexpanded{,typepedagolibelle=}\cmdseqca@planche@typepedagolibelle%
		\unexpanded{,codecouleur=}\cmdseqca@planche@codecouleur \unexpanded{]}%
		\@charlb \unexpanded{#1} \@charrb
		% Séparateur : \\ toutes les 4 cellules, & sinon
		\ifnum\numexpr\value{seqca@cellpos}-(\value{seqca@cellpos}/4)*4\relax=0 %
			\unexpanded{\\}%
		\else
			\unexpanded{&}%
		\fi
	}%
}%
% \end{macrocode}%
% \end{macro}%
% -----------------------------------------------------------------------------
% Ecrit le verso dans le fichier annexe
% -----------------------------------------------------------------------------
% \begin{macro}{\seqca@emettreVerso}%
% #1 : contenu LaTeX du verso
% \begin{macrocode}%
\newcommand{\seqca@emettreVerso}[1]{%
	\stepcounter{seqca@cellpos}%
	\immediate\write\plancheversostream{%
		\unexpanded{\seqCarteVerso[niveau=}\cmdseqca@planche@niveau%
		\unexpanded{,sequence=}\cmdseqca@planche@sequence%
		\unexpanded{,num=}\cmdseqca@planche@num%
		\unexpanded{,codecouleur=}\cmdseqca@planche@codecouleur \unexpanded{]}%
		\@charlb \unexpanded{#1} \@charrb
		\ifnum\numexpr\value{seqca@cellpos}-(\value{seqca@cellpos}/4)*4\relax=0 %
			\unexpanded{\\}%
		\else
			\unexpanded{&}%
		\fi
	}%
}%
% \end{macrocode}%
% \end{macro}%

% =============================================================================
% MACROS ET ENVIRONNEMENTS PUBLICS
% =============================================================================
%
% -----------------------------------------------------------------------------
% Application d'un thème de couleurs
% -----------------------------------------------------------------------------
\newcommand{\seqThemeCarte}[1]{
	\seq@setColors{#1}
	\seqca@applytheme
}
% -----------------------------------------------------------------------------
% Initialisations pour les planches de cartes paramétrées
% -----------------------------------------------------------------------------
% \begin{macro}{\seqInitCarteParam}%
% \IndexMacro{\seqInitCarteParam}%
% \begin{macrocode}%
% Désactive les shorthands de babel+french pour xintexpr,
% rend robustes les macros qui risquent de casser dans \write
% et crée les flux d'écriture pour les planches et 
% pour les recap
% A appeler avant \seqPlanche
\newcommand\seqInitCartes{%
	\shorthandoff{:;?!}
	\newwrite\plancherectostream
	\newwrite\plancheversostream
	\newwrite\tiragesstream
	\newwrite\recaprectostream
	\newwrite\recapversostream
%	\robustify{\smallskip}
%	\robustify{\medskip}
%	\robustify{\bigskip}
%	\robustify{\seqFrac}
%	\robustify{\textsuperscript}
%	\robustify{\textsubscript}
%	\robustify{\newline}
%	\robustify{\textbf}
%	\robustify{\text}
%	\robustify{\degre}
}%
% \end{macrocode}%
% \end{macro}%
% -----------------------------------------------------------------------------
% Composition du recto
% -----------------------------------------------------------------------------
% \begin{macro}{\seqCarteRecto}%
% \IndexMacro{\seqCarteRecto}%
% #1 : options
% #2 : contenu LaTeX du recto (question)
% \begin{macrocode}%
\define@cmdkey[seqca]{recto}{niveau}[N10]{}%
\define@cmdkey[seqca]{recto}{sequence}[S01]{}%
\define@cmdkey[seqca]{recto}{num}[1]{}%
\define@cmdkey[seqca]{recto}{typepedagolibelle}[Définition]{}%
\define@cmdkey[seqca]{recto}{codecouleur}[nombres]{}%
\newcommand{\seqCarteRecto}[2][]{%
    \setkeys[seqca]{recto}{niveau,sequence,num,typepedagolibelle,%
                           codecouleur,#1}%
    \seq@setColors{\cmdseqca@recto@codecouleur}%
    \seqca@applytheme
    \seqca@recto
        {\cmdseqca@recto@niveau}%
        {\cmdseqca@recto@sequence}%
        {\cmdseqca@recto@num}%
        {\cmdseqca@recto@typepedagolibelle}%
        {#2}%
		\ignorespaces
}%
% \end{macrocode}%
% \end{macro}%
% -----------------------------------------------------------------------------
% Composition du verso
% -----------------------------------------------------------------------------
% \begin{macro}{\seqCarteVerso}%
% \IndexMacro{\seqCarteVerso}%
% #1 : options
% #2 : contenu LaTeX du verso (réponse)
% \begin{macrocode}%
\define@cmdkey[seqca]{verso}{niveau}[N10]{}%
\define@cmdkey[seqca]{verso}{sequence}[S01]{}%
\define@cmdkey[seqca]{verso}{num}[1]{}%
\define@cmdkey[seqca]{verso}{codecouleur}[nombres]{}%
\newcommand{\seqCarteVerso}[2][]{%
    \setkeys[seqca]{verso}{niveau,sequence,num,%
                           codecouleur,#1}%
    \seq@setColors{\cmdseqca@verso@codecouleur}%
    \seqca@applytheme
    \seqca@verso
        {\cmdseqca@verso@niveau}%
        {\cmdseqca@verso@sequence}%
        {\cmdseqca@verso@num}%
        {#2}%
        \ignorespaces
}%
% \end{macrocode}%
% \end{macro}%
% -----------------------------------------------------------------------------
% Composition d'une carte d'automatisme complète
% -----------------------------------------------------------------------------
% \begin{macro}{\seqCarteAutomatisme}%
% \IndexMacro{\seqCarteAuto}%
% Produit une carte d'automatisme : 1 recto (question) + 1 verso (réponse).
% Si la carte est paramétrée, les paramètres doivent être insérés dans le
% code avant l'appel à cette macro
% #1 : options
% #2 : contenu LaTeX du recto
% #3 : contenu LaTeX du verso
% \begin{macrocode}%
%% Options
\define@cmdkey[seqca]{carte}{niveau}[N10]{}%
\define@cmdkey[seqca]{carte}{sequence}[S01]{}%
\define@cmdkey[seqca]{carte}{num}[1]{}%
\define@cmdkey[seqca]{carte}{typepedagolibelle}[Définition]{}%
\define@cmdkey[seqca]{carte}{codecouleur}[nombres]{}%
\define@cmdkey[seqca]{carte}{nom}[]{}%
\newcommand{\seqCarteAutomatisme}[3][]{%
    \setkeys[seqca]{carte}{niveau,sequence,num,typepedagolibelle,%
                           codecouleur,nom,#1}%
    \seq@setColors{\cmdseqca@carte@codecouleur}%
    \seqca@applytheme
    % Recto
    \seqca@recto
        {\cmdseqca@carte@niveau}%
        {\cmdseqca@carte@sequence}%
        {\cmdseqca@carte@num}%
        {\cmdseqca@carte@typepedagolibelle}%
        {#2}%
    % Verso
    \seqca@verso
        {\cmdseqca@carte@niveau}%
        {\cmdseqca@carte@sequence}%
        {\cmdseqca@carte@num}%
        {#3}%
}%
% \end{macrocode}%
% \end{macro}%
% -----------------------------------------------------------------------------
% Initialisation d'une double page de récap
% -----------------------------------------------------------------------------
%\begin{environment}{seqCarteRecap}%
% \begin{environment}{seqCarteRecap}%
% \IndexEnv{seqCartePageRecap}%
% Jusqu'à 4 lignes de 4 cartes, rectos sur la 1ere page, versos sur la 2nde.
% \begin{macrocode}%
%% Environnement
\newenvironment{seqCarteRecap}{%
	\immediate\openout\recaprectostream=recaprecto.tex\relax
	\immediate\openout\recapversostream=recapverso.tex\relax
	\immediate\write\recaprectostream{\unexpanded{\begin{tblr}{*{4}{X[c]}}}}
	\immediate\write\recapversostream{\unexpanded{\begin{tblr}{*{4}{X[c]}}}}
}{
	\immediate\write\recaprectostream{\unexpanded{\end{tblr}\clearpage}}
	\immediate\write\recapversostream{\unexpanded{\end{tblr}\clearpage}}
	\immediate\closeout\recaprectostream
	\immediate\closeout\recapversostream
	\input{recaprecto}%
	\input{recapverso}%
}
% \end{macrocode}%
% \end{environment}%
% -----------------------------------------------------------------------------
% Ajout d'une carte d'automatisme complète à un récap
% -----------------------------------------------------------------------------
% \begin{macro}{\seqCarteRecapAjouteCarte}%
% \IndexMacro{\seqCarteRecapAjouteCarte}%
% 
% #1 : options
% #2 : variables
% #3 : contenu LaTeX du recto
% #4 : contenu LaTeX du verso
% \begin{macrocode}%
%% Options
\define@cmdkey[seqca]{addrecap}{niveau}[N10]{}%
\define@cmdkey[seqca]{addrecap}{sequence}[S01]{}%
\define@cmdkey[seqca]{addrecap}{num}[1]{}%
\define@cmdkey[seqca]{addrecap}{typepedagolibelle}[Définition]{}%
\define@cmdkey[seqca]{addrecap}{codecouleur}[nombres]{}%
\define@cmdkey[seqca]{addrecap}{finLigne}[non]{}%
%% Macro
\newcommand{\seqCarteRecapAjouteCarte}[4][]{
	\setkeys[seqca]{addrecap}{niveau,sequence,num,typepedagolibelle,%
			codecouleur,finLigne,#1}%
	% Variables
	#2
	% Recto
	\immediate\write\recaprectostream{\unexpanded{\seqCarteRecto[niveau=}\cmdseqca@addrecap@niveau%
			\unexpanded{,sequence=}\cmdseqca@addrecap@sequence%
			\unexpanded{,num=}\cmdseqca@addrecap@num%
			\unexpanded{,typepedagolibelle=}\cmdseqca@addrecap@typepedagolibelle%
			\unexpanded{,codecouleur=}\cmdseqca@addrecap@codecouleur \unexpanded{]}%
			\@charlb \unexpanded{#3} \@charrb%
	}%
	\ifdefstring{\cmdseqca@addrecap@finLigne}{oui}{
		\immediate\write\recaprectostream{\\}
	}{
		\immediate\write\recaprectostream{&}
	}
	% Verso
	\immediate\write\recapversostream{\unexpanded{\seqCarteVerso[niveau=}\cmdseqca@addrecap@niveau%
			\unexpanded{,sequence=}\cmdseqca@addrecap@sequence%
			\unexpanded{,num=}\cmdseqca@addrecap@num%
			\unexpanded{,codecouleur=}\cmdseqca@addrecap@codecouleur \unexpanded{]}
			\@charlb \unexpanded{#4} \@charrb%
	}%
	\ifdefstring{\cmdseqca@addrecap@finLigne}{oui}{
		\immediate\write\recapversostream{\\}
	}{
		\immediate\write\recapversostream{&}
	}
}
% \end{macrocode}%
% \end{macro}%
% -----------------------------------------------------------------------------
% Ajout une ligne de 4 cartes d'automatisme complètes à un récap
% -----------------------------------------------------------------------------
% \begin{macro}{\seqCarteRecapAjouteLigneCartes}%
% \IndexMacro{\seqCarteRecapAjouteLigneCartes}%
% 
% #1 : options
% #2 : variables
% \begin{macrocode}%
%% Options
\define@cmdkey[seqca]{addlignerecap}{niveau}[N10]{}%
\define@cmdkey[seqca]{addlignerecap}{sequenceUn}[S01]{}%
\define@cmdkey[seqca]{addlignerecap}{numUn}[1]{}%
\define@cmdkey[seqca]{addlignerecap}{typepedagolibelleUn}[Définition]{}%
\define@cmdkey[seqca]{addlignerecap}{codecouleurUn}[nombres]{}%
\define@key[seqca]{addlignerecap}{rectoUn}[test_recto_un]{\long\def\cmdseqca@addlignerecap@rectoUn{#1}}%
\define@key[seqca]{addlignerecap}{versoUn}[test_verso_un]{\long\def\cmdseqca@addlignerecap@versoUn{#1}}%
%\define@cmdkey[seqca]{addlignerecap}{rectoUn}[test_recto_un]{}%
%\define@cmdkey[seqca]{addlignerecap}{versoUn}[test_verso_un]{}%
\define@cmdkey[seqca]{addlignerecap}{sequenceDeux}[S01]{}%
\define@cmdkey[seqca]{addlignerecap}{numDeux}[2]{}%
\define@cmdkey[seqca]{addlignerecap}{typepedagolibelleDeux}[Définition]{}%
\define@cmdkey[seqca]{addlignerecap}{codecouleurDeux}[nombres]{}%
\define@key[seqca]{addlignerecap}{rectoDeux}[test_recto_deux]{\long\def\cmdseqca@addlignerecap@rectoDeux{#1}}%
\define@key[seqca]{addlignerecap}{versoDeux}[test_verso_deux]{\long\def\cmdseqca@addlignerecap@versoDeux{#1}}%
%\define@cmdkey[seqca]{addlignerecap}{rectoDeux}[test_recto_deux]{}%
%\define@cmdkey[seqca]{addlignerecap}{versoDeux}[test_verso_deux]{}%
\define@cmdkey[seqca]{addlignerecap}{sequenceTrois}[S01]{}%
\define@cmdkey[seqca]{addlignerecap}{numTrois}[3]{}%
\define@cmdkey[seqca]{addlignerecap}{typepedagolibelleTrois}[Définition]{}%
\define@cmdkey[seqca]{addlignerecap}{codecouleurTrois}[nombres]{}%
\define@key[seqca]{addlignerecap}{rectoTrois}[test_recto_trois]{\long\def\cmdseqca@addlignerecap@rectoTrois{#1}}%
\define@key[seqca]{addlignerecap}{versoTrois}[test_verso_trois]{\long\def\cmdseqca@addlignerecap@versoTrois{#1}}%
%\define@cmdkey[seqca]{addlignerecap}{rectoTrois}[test_recto_trois]{}%
%\define@cmdkey[seqca]{addlignerecap}{versoTrois}[test_verso_trois]{}%
\define@cmdkey[seqca]{addlignerecap}{sequenceQuatre}[S01]{}%
\define@cmdkey[seqca]{addlignerecap}{numQuatre}[4]{}%
\define@cmdkey[seqca]{addlignerecap}{typepedagolibelleQuatre}[Définition]{}%
\define@cmdkey[seqca]{addlignerecap}{codecouleurQuatre}[nombres]{}%
\define@key[seqca]{addlignerecap}{rectoQuatre}[test_recto_quatre]{\long\def\cmdseqca@addlignerecap@rectoQuatre{#1}}%
\define@key[seqca]{addlignerecap}{versoQuatre}[test_verso_quatre]{\long\def\cmdseqca@addlignerecap@versoQuatre{#1}}%
%\define@cmdkey[seqca]{addlignerecap}{rectoQuatre}[test_recto_quatre]{}%
%\define@cmdkey[seqca]{addlignerecap}{versoQuatre}[test_verso_quatre]{}%
%% Macro
\NewDocumentCommand{\seqCarteRecapAjouteLigneCartes}{+O{} +m}{
%\newcommand{\seqCarteRecapAjouteLigneCartes}[2][]{
	\setkeys[seqca]{addlignerecap}{niveau,%
			sequenceUn,numUn,typepedagolibelleUn,codecouleurUn,rectoUn,versoUn,%
			sequenceDeux,numDeux,typepedagolibelleDeux,codecouleurDeux,rectoDeux,versoDeux,%
			sequenceTrois,numTrois,typepedagolibelleTrois,codecouleurTrois,rectoTrois,versoTrois,%
			sequenceQuatre,numQuatre,typepedagolibelleQuatre,codecouleurQuatre,rectoQuatre,versoQuatre,%
			#1}%
	% Variables
	#2
	% Rectos
	\immediate\write\recaprectostream{\unexpanded{\seqCarteRecto[niveau=}\cmdseqca@addlignerecap@niveau%
			\unexpanded{,sequence=}\cmdseqca@addlignerecap@sequenceUn%
			\unexpanded{,num=}\cmdseqca@addlignerecap@numUn%
			\unexpanded{,typepedagolibelle=}\cmdseqca@addlignerecap@typepedagolibelleUn%
			\unexpanded{,codecouleur=}\cmdseqca@addlignerecap@codecouleurUn \unexpanded{]}%
			\@charlb \unexpanded\expandafter{\cmdseqca@addlignerecap@rectoUn} \@charrb &%
	}%
	\immediate\write\recaprectostream{\unexpanded{\seqCarteRecto[niveau=}\cmdseqca@addlignerecap@niveau%
			\unexpanded{,sequence=}\cmdseqca@addlignerecap@sequenceDeux%
			\unexpanded{,num=}\cmdseqca@addlignerecap@numDeux%
			\unexpanded{,typepedagolibelle=}\cmdseqca@addlignerecap@typepedagolibelleDeux%
			\unexpanded{,codecouleur=}\cmdseqca@addlignerecap@codecouleurDeux \unexpanded{]}%
			\@charlb \unexpanded\expandafter{\cmdseqca@addlignerecap@rectoDeux} \@charrb &%
	}%
	\immediate\write\recaprectostream{\unexpanded{\seqCarteRecto[niveau=}\cmdseqca@addlignerecap@niveau%
			\unexpanded{,sequence=}\cmdseqca@addlignerecap@sequenceTrois%
			\unexpanded{,num=}\cmdseqca@addlignerecap@numTrois%
			\unexpanded{,typepedagolibelle=}\cmdseqca@addlignerecap@typepedagolibelleTrois%
			\unexpanded{,codecouleur=}\cmdseqca@addlignerecap@codecouleurTrois \unexpanded{]}%
			\@charlb \unexpanded\expandafter{\cmdseqca@addlignerecap@rectoTrois} \@charrb &%
	}%
	\immediate\write\recaprectostream{\unexpanded{\seqCarteRecto[niveau=}\cmdseqca@addlignerecap@niveau%
			\unexpanded{,sequence=}\cmdseqca@addlignerecap@sequenceQuatre%
			\unexpanded{,num=}\cmdseqca@addlignerecap@numQuatre%
			\unexpanded{,typepedagolibelle=}\cmdseqca@addlignerecap@typepedagolibelleQuatre%
			\unexpanded{,codecouleur=}\cmdseqca@addlignerecap@codecouleurQuatre \unexpanded{]}%
			\@charlb \unexpanded\expandafter{\cmdseqca@addlignerecap@rectoQuatre} \@charrb \\%
	}%
	% Verso
	\immediate\write\recapversostream{\unexpanded{\seqCarteVerso[niveau=}\cmdseqca@addlignerecap@niveau%
			\unexpanded{,sequence=}\cmdseqca@addlignerecap@sequenceQuatre%
			\unexpanded{,num=}\cmdseqca@addlignerecap@numQuatre%
			\unexpanded{,codecouleur=}\cmdseqca@addlignerecap@codecouleurQuatre \unexpanded{]}
			\@charlb \unexpanded\expandafter{\cmdseqca@addlignerecap@versoQuatre} \@charrb &%
	}%
	\immediate\write\recapversostream{\unexpanded{\seqCarteVerso[niveau=}\cmdseqca@addlignerecap@niveau%
			\unexpanded{,sequence=}\cmdseqca@addlignerecap@sequenceTrois%
			\unexpanded{,num=}\cmdseqca@addlignerecap@numTrois%
			\unexpanded{,codecouleur=}\cmdseqca@addlignerecap@codecouleurTrois \unexpanded{]}
			\@charlb \unexpanded\expandafter{\cmdseqca@addlignerecap@versoTrois} \@charrb &%
	}%
	\immediate\write\recapversostream{\unexpanded{\seqCarteVerso[niveau=}\cmdseqca@addlignerecap@niveau%
			\unexpanded{,sequence=}\cmdseqca@addlignerecap@sequenceDeux%
			\unexpanded{,num=}\cmdseqca@addlignerecap@numDeux%
			\unexpanded{,codecouleur=}\cmdseqca@addlignerecap@codecouleurDeux \unexpanded{]}
			\@charlb \unexpanded\expandafter{\cmdseqca@addlignerecap@versoDeux} \@charrb &%
	}%
	\immediate\write\recapversostream{\unexpanded{\seqCarteVerso[niveau=}\cmdseqca@addlignerecap@niveau%
			\unexpanded{,sequence=}\cmdseqca@addlignerecap@sequenceUn%
			\unexpanded{,num=}\cmdseqca@addlignerecap@numUn%
			\unexpanded{,codecouleur=}\cmdseqca@addlignerecap@codecouleurUn \unexpanded{]}
			\@charlb \unexpanded\expandafter{\cmdseqca@addlignerecap@versoUn} \@charrb \\%
	}%
}
% \end{macrocode}%
% \end{macro}%
% -----------------------------------------------------------------------------
% Composition d'une planche complète d'une carte (fixe ou paramétrée)
% -----------------------------------------------------------------------------
% \begin{macro}{\seqCartePlanche}%
% \IndexMacro{\seqCartePlanche}%
% 
% #1 : options
% #2 : variables (optionnel)
% #3 : contenu LaTeX du recto
% #4 : contenu LaTeX du verso
% \begin{macrocode}%
\define@cmdkey[seqca]{planche}{niveau}[N10]{}%
\define@cmdkey[seqca]{planche}{sequence}[S01]{}%
\define@cmdkey[seqca]{planche}{num}[1]{}%
\define@cmdkey[seqca]{planche}{typepedagolibelle}[Définition]{}%
\define@cmdkey[seqca]{planche}{codecouleur}[nombres]{}%
\newcommand{\seqCartePlanche}[4][]{
	\setkeys[seqca]{planche}{niveau,sequence,num,typepedagolibelle,%
			codecouleur,#1}%
	\immediate\openout\plancherectostream=plancherecto.tex\relax
	\immediate\openout\plancheversostream=plancheverso.tex\relax
	\immediate\openout\tiragesstream=tirages.tex\relax
	% Ecrire les tirages figés dans le fichier annexe
	\immediate\write\tiragesstream{\unexpanded{#2}}%
	% Ecrire l'ouverture des tableaux
	\immediate\write\plancherectostream{\unexpanded{\begin{tblr}{*{4}{X[c]}}}}%
	\immediate\write\plancheversostream{\unexpanded{\begin{tblr}{*{4}{X[c]}}}}%
	% Ecrire les recto 
	\begingroup
		\setcounter{seqca@cellpos}{0}%
		\let\seqCelluleRecto\seqca@emettreRecto
		#3%
	\endgroup
	% Ecrire les verso
	\begingroup
		\setcounter{seqca@cellpos}{0}%
		\let\seqCelluleVerso\seqca@emettreVerso
		#4%
	\endgroup
	% Ecrire la fermeture des tableaux et les sauts de page
	\immediate\write\plancherectostream{\unexpanded{\end{tblr}\clearpage}}%
	\immediate\write\plancheversostream{\unexpanded{\end{tblr}\clearpage}}%
	% Fermer les fichiers annexes
	\immediate\closeout\tiragesstream\relax
	\immediate\closeout\plancherectostream\relax
	\immediate\closeout\plancheversostream\relax
	% Inclure les fichiers
	\input{tirages}%
	\input{plancherecto}%
	\input{plancheverso}%
}
% \end{macrocode}%
% \end{macro}%

% Fin du module carte-automatisme
%</package>
